% =============================================================================
% 6.5 LÓGICA PROGRAMABLE EN EL MOTOR DE BASE DE DATOS
% =============================================================================
\section{Lógica Programable en el Motor de Base de Datos}
\label{sec:programable}

EthosHub traslada una parte significativa de la lógica de negocio sensible al propio
motor PostgreSQL, mediante \textbf{\textit{triggers}} (disparadores) y
\textbf{funciones PL/pgSQL}. Esta decisión persigue dos objetivos: garantizar
invariantes que no deben depender de la aplicación (auditoría, exclusividad) y operar de
forma segura por encima de las políticas RLS mediante funciones \texttt{SECURITY DEFINER}.

% -----------------------------------------------------------------------------
\subsection{Triggers para auditoría y sincronización de eventos}
\label{subsec:triggers}

El esquema define 12 \textit{triggers}, clasificables en cuatro familias funcionales,
resumidas en la tabla~\ref{tab:triggers}.

\clearpage
\vspace{0.2cm}
\noindent
\renewcommand{\arraystretch}{1.25}
\fontsize{8.2}{9.8}\selectfont\sffamily
\begin{tabularx}{\textwidth}{|>{\ttfamily\color{colorPrimario}\scriptsize}l|>{\scriptsize}l|X|}
\hline
\rowcolor{headerTabla}
\sffamily\textbf{Trigger} & \sffamily\textbf{Tabla / Evento} & \sffamily\textbf{Función} \\
\hline
notify\_chat\_message & chat\_messages / AFTER INSERT & Crea una notificación para el destinatario del mensaje. \\
\hline
\rowcolor{grisTecnico}
trg\_chat\_message\_timestamp & chat\_messages / AFTER INSERT & Actualiza \texttt{updated\_at} del chat (orden de conversaciones). \\
\hline
notify\_profile\_like & profile\_likes / AFTER INSERT & Notifica al candidato que ha recibido un \textit{like}. \\
\hline
\rowcolor{grisTecnico}
trg\_cv\_documents\_updated\_at & cv\_documents / BEFORE UPDATE & Mantiene \texttt{updated\_at} en cada modificación. \\
\hline
trg\_profile\_connections\_updated\_at & profile\_connections / BEFORE UPDATE & Actualiza \texttt{updated\_at}. \\
\hline
\rowcolor{grisTecnico}
trg\_check\_basic\_exclusivity & profiles\_basic / BEFORE INSERT & Impide que un usuario sea profesional y reclutador a la vez. \\
\hline
trg\_check\_company\_exclusivity & profiles\_company / BEFORE INSERT & Misma exclusividad para el perfil reclutador. \\
\hline
\rowcolor{grisTecnico}
log\_profile\_basic\_events & profiles\_basic / INSERT,UPDATE & Audita altas y borrados en \texttt{admin.system\_logs}. \\
\hline
log\_project\_created & projects / AFTER INSERT & Registro de auditoría de creación de proyectos. \\
\hline
\rowcolor{grisTecnico}
log\_portfolio\_visibility & portfolio\_settings / UPDATE & Audita cambios de publicación del portafolio. \\
\hline
log\_login\_blocked & login\_attempts / UPDATE & Audita bloqueos por fuerza bruta. \\
\hline
\rowcolor{grisTecnico}
log\_skill\_tag\_created & global\_skill\_tags / AFTER INSERT & Audita la creación de nuevas etiquetas de habilidad. \\
\hline
\end{tabularx}
\label{tab:triggers}

El siguiente \textit{trigger} ejemplifica la \textbf{sincronización de eventos}: al
insertarse un mensaje, genera automáticamente una notificación para el destinatario,
resolviendo quién es el receptor a partir de la conversación.

\begin{lstlisting}[language=SQL, caption={Trigger de notificación de mensaje (\texttt{core.notify\_chat\_message}).}]
CREATE FUNCTION core.notify_chat_message() RETURNS trigger
  LANGUAGE plpgsql SECURITY DEFINER SET search_path TO 'core','pg_temp' AS $$
DECLARE v_recipient UUID;
BEGIN
  SELECT CASE WHEN c.basic_profile_id = NEW.sender_id THEN c.company_profile_id
              ELSE c.basic_profile_id END
  INTO v_recipient FROM core.chats c WHERE c.chat_id = NEW.chat_id;

  IF v_recipient IS NOT NULL AND v_recipient <> NEW.sender_id THEN
    INSERT INTO core.notifications (profile_id, type, title, message)
    VALUES (v_recipient, 'message', 'Nuevo mensaje',
            left(COALESCE(NULLIF(NEW.content,''), 'Has recibido un adjunto'), 120));
  END IF;
  RETURN NEW;
END; $$;
\end{lstlisting}

\paragraph{Auditoría histórica SCD Tipo 2.} El esquema \comando{history} implementa
\textbf{Slowly Changing Dimensions de Tipo 2} sobre los perfiles: cada cambio relevante
genera una nueva versión histórica (con marcas de vigencia), preservando el rastro
completo de la evolución de cada perfil en \texttt{history.hist\_profiles} y
\texttt{history.hist\_profiles\_basic}.

% -----------------------------------------------------------------------------
\subsection{Procedimientos almacenados y funciones PL/pgSQL}
\label{subsec:funciones}

El esquema expone alrededor de \textbf{60 funciones}, la mayoría declaradas como
\comando{SECURITY DEFINER} para operar de forma controlada por encima de RLS. Se
agrupan funcionalmente como muestra la tabla~\ref{tab:rpcs}.

\vspace{0.2cm}
\noindent
\renewcommand{\arraystretch}{1.25}
\fontsize{8.2}{9.8}\selectfont\sffamily
\begin{tabularx}{\textwidth}{|>{\bfseries\color{colorPrimario}}l|X|}
\hline
\rowcolor{headerTabla}
\textbf{Grupo} & \textbf{Funciones representativas (RPC)} \\
\hline
Chat & \texttt{send\_message}, \texttt{get\_messages}, \texttt{mark\_messages\_read}, \texttt{upsert\_chat}, \texttt{delete\_chat}, \texttt{get\_chat\_list}, \texttt{get\_unread\_count}. \\
\hline
\rowcolor{grisTecnico}
Pipeline / Likes & \texttt{add\_like}, \texttt{remove\_like}, \texttt{update\_like\_stage}, \texttt{get\_company\_likes}, \texttt{has\_liked\_profile}. \\
\hline
Talento & \texttt{record\_talent\_view}, \texttt{get\_recent\_talent\_views}, \texttt{get\_explore\_portfolios}, \texttt{get\_profile\_likes\_count}. \\
\hline
\rowcolor{grisTecnico}
Seguridad & \texttt{check\_login\_block}, \texttt{record\_login\_failure}, \texttt{clear\_login\_attempts}, \texttt{increment\_otp\_attempt}, \texttt{cleanup\_expired\_codes}. \\
\hline
Identidad & \texttt{handle\_new\_user}, \texttt{delete\_profile\_cascade}, \texttt{sync\_auth\_email\_to\_profiles}, \texttt{sync\_oauth\_identity\_to\_connections}. \\
\hline
\rowcolor{grisTecnico}
Administración & \texttt{admin.search\_profiles}, \texttt{admin.get\_global\_metrics}, \texttt{admin.soft\_delete\_profile}, \texttt{admin.restore\_profile}, \texttt{admin.log\_event}. \\
\hline
\end{tabularx}
\label{tab:rpcs}

La función \comando{record\_login\_failure} ilustra la lógica anti fuerza bruta
ejecutada de forma atómica en el motor: incrementa el contador, reinicia si el bloqueo
expiró y aplica un bloqueo de 15 minutos al quinto fallo, todo en una sola sentencia
\texttt{UPSERT}.

\begin{lstlisting}[language=SQL, caption={Función anti fuerza bruta (\texttt{core.record\_login\_failure}, fragmento).}]
INSERT INTO core.login_attempts (email, failed_count, last_attempt)
VALUES (lower(p_email), 1, now())
ON CONFLICT (lower(email)) DO UPDATE
  SET failed_count = CASE
        WHEN login_attempts.blocked_until < now() THEN 1          -- bloqueo expirado: reinicia
        ELSE login_attempts.failed_count + 1 END,
      blocked_until = CASE
        WHEN (...failed_count + 1) >= 5                            -- 5 intentos
        THEN now() + '15 minutes'::INTERVAL ELSE NULL END,
      last_attempt = now()
RETURNING *;
\end{lstlisting}

\begin{AlertaSeguridad}
El uso de \comando{SECURITY DEFINER} es potente pero delicado: estas funciones se
ejecutan con los privilegios de su propietario, eludiendo RLS. Por ello \textbf{fijan
explícitamente} \comando{SET search\_path} (p.~ej. \texttt{'core','pg\_temp'}) para
prevenir ataques de secuestro de \textit{search path}, y validan internamente la
pertenencia del llamante antes de operar sobre datos ajenos.
\end{AlertaSeguridad}
